\parttitle{XI}{From Pipeline to Product: Extreme-Climate Forecasting for Users}

Part~X delivers a validated forecast core. This part wraps it in the full
system a user would touch --- data plane, hybrid forecast core, decision
layer, product surfaces, personas --- adds implementation Phases~8--12, and
integrates a survey of classical ML/DL and agentic-AI post-processing at the
user boundary. The governance is inherited unchanged: pre-registered
promotion gates, simulation-scoped claims, every number traceable to a run
identifier, and the quantum contribution always reported as a delta against
its own classical ablation.

\section{Product definition}
\label{sec:prod-def}

\emph{Two tracks, routed by measured memory requirements.} Track~A
(operational extremes, minutes to hours): solar irradiance and PV-output
ramp events driven by convective cloud fields, squalls, and haze --- the
extreme-climate signal at operational timescale, and the regime where the
evidence of Part~X locates QRC value; short-memory, high-nonlinearity
targets, served by RF-QRC with virtual nodes. Flagship geography: tropical
convective regimes (Singapore and South-East Asian rooftop and microgrid
PV), where ramps are frequent, steep, and costly. Track~B (tactical and
seasonal extremes, weeks to months): ENSO phase transitions and
anomalous-heat episodes; long-nonlinear-memory targets, served by the
recurrent CZ-ring reservoir whose ENSO advantage this handbook measured
(Part~VII). The Part~X ``ENSO tension'' thereby becomes an
architecture-selection rule rather than an embarrassment.

\emph{Decisions fed and error costs, in writing.} Track~A feeds battery
pre-charge and discharge, spinning-reserve commitment, curtailment, and
intraday positions; a miss costs imbalance penalties, diesel spin-up, or
voltage excursions, while a false alarm costs battery cycling and reserve
over-procurement. Track~B feeds seasonal hedging, maintenance windows,
water and agricultural allocation, and parametric-insurance triggers.

\emph{Incumbents named.} Persistence and smart-persistence (the honest
one-step champion on solar --- this project's own calibration); NWP-derived
vendor forecasts; all-sky-imager nowcast systems, whose published ramp
capture spans roughly $26$--$92\%$ with large false-alarm spreads --- a
genuinely weak incumbent frontier on ramps; and deep-learning weather
models, which the product \emph{consumes} as exogenous channels rather than
competes with. The venture owns three things only: the event-focused hybrid
feature layer, the calibrated decision layer that converts forecasts into
cost-optimal alerts, and the benchmark-governance layer --- the
no-overclaiming discipline packaged as a feature. The quantum module sits
behind an interface: if its pre-registered delta over the classical-only
ablation fails at any promotion review, the product ships classical-only
and the quantum track continues as research. Quantum-optional architecture
is the de-risking posture.

\emph{Why RF-QRC specifically is the productisable QRC.} Recurrent QRC
carries quantum state across steps: inference is stateful, sequential,
non-cacheable, and latency-bound to the quantum backend. RF-QRC step
features depend only on the current quantised input window; memory is the
classical leak. Features are therefore memoisable, per-step computation is
stateless and parallel, and inference latency decouples from the backend.
At product-relevant sizes ($n\le11$) the exact simulator \emph{is} the
production backend --- fast, deterministic, and identical in distribution
to ideal hardware --- while the reuse compiler keeps the hardware path open
for the regime where simulation fails, at compiled widths independent of
memory horizon (Section~\ref{sec:rfqrc-reuse-fit}). This is stated plainly
to users: near-term inference is quantum-inspired simulation; hardware
execution is a scaling option, not a marketing prop.

\subsection{Personas and their contracts}
\label{sec:prod-personas}

Every persona contract names its metric before onboarding --- the
pre-registration discipline applied to customers.

\begin{table*}[h]
\centering\footnotesize
\caption{Personas, decisions, and the metric each contract names.}
\label{tab:prod-personas}
\begin{tabular}{@{}p{2.9cm}p{3.1cm}p{1.7cm}p{2.6cm}p{3.4cm}@{}}
\toprule
Persona & Decision & Horizon & Surface & Metric that matters \\
\midrule
microgrid / island-grid operator & battery and genset scheduling & 15\,min--4\,h & alerts (chat/SMS/webhook), dashboard & ramp F$_1$ at a fixed false-alarm budget; lead time \\
C\&I rooftop energy manager & load shifting, demand-charge avoidance & 30\,min--6\,h & dashboard, weekly digest & avoided cost per month \\
utility control room / VPP & reserve procurement, curtailment & 1--24\,h & API into EMS/SCADA historian & pinball loss on quantiles; recall at capped false-alarm rate \\
energy trader & intraday position vs imbalance & 1--12\,h & API, scenario tool & cost-weighted CRPS; volatility-regime hit rate \\
climate-risk analyst / insurer & parametric triggers, seasonal exposure & weeks--months & reports, API & calibrated exceedance probabilities \\
researcher tier & benchmarking, method development & --- & open harness, feature API & seed-exact reproducibility \\
\bottomrule
\end{tabular}
\end{table*}

\section{End-to-end architecture}
\label{sec:prod-arch}

\begin{strip}
\begin{lstlisting}[language={},basicstyle=\ttfamily\scriptsize,caption={Six-layer architecture. L6 governance spans the whole pipeline.},label={lst:prod-layers}]
L1 INGEST      L2 TRANSFORM        L3 FORECAST CORE       L4 DECISION       L5 SURFACES
SURFRAD/NSRDB  QC (mask, never     classical features     calibration       API (REST/stream)
ERA5 / NOAA    interpolate)        + RF-QRC+VN features   (conformal,       dashboard
site SCADA  -> clear-sky index  -> via FeatureProvider -> isotonic)      -> alert engine
NWP + DL-      k_t (train-span     [sim|hw-batch|cache]   stacker (GBM/     agent assistant
weather feeds  fit, frozen per     quantile+event heads   shallow DL)
(exogenous)    version); anomaly   shadow baseline        regime router
               transforms;         battery (always on)    cost-sensitive
               feature store                              thresholds
        L6 GOVERNANCE: drift + calibration monitors, effective-rank health,
        champion-challenger, promotion gates = pre-registered criteria,
        model cards, run-ID audit log
\end{lstlisting}
\end{strip}

\emph{L1, ingestion and QC.} Connectors per the dataset registry of
Part~VI plus private site feeds; the vetting checklist applies to every new
source; missing data is masked, never silently interpolated; gaps longer
than the autocorrelation time split the series; raw data is immutable and
versioned.

\emph{L2, transform and feature store.} Physics-derived target transforms
(solar to clear-sky index; climate to anomalies against training-years
climatology) live in a transform service holding the frozen train-span
statistics per model version --- the leakage firewall as an architectural
component rather than a code-review hope. Calendar and solar-geometry
features come from lookup: no quantum resource is spent on what a lookup
table solves.

\emph{L3, forecast core.} The Part~X core behind a narrow interface:

\begin{strip}
\begin{lstlisting}[language=Python,caption={The forecast core's only coupling to quantum execution.},label={lst:prod-provider}]
class FeatureProvider(Protocol):
    def features(self, window: np.ndarray, cfg: RFQRCConfig) -> np.ndarray: ...

# SimProvider      exact Aer snapshots (production default, n <= 11)
# HardwareProvider reuse-compiled batched jobs (qreuse_batch -> compile_with_reuse),
#                  gated by the self-calibrated TVD equivalence check as CI
# CachedProvider   memo keyed by (config_hash, quantised window m); valid because
#                  RF-QRC step features are stateless; leak applied downstream
\end{lstlisting}
\end{strip}

Heads on the shared feature stream: multi-horizon quantile read-outs
trained on pinball loss (one $W_{\mathrm{out}}$ per horizon), an
event-probability head, and a per-horizon leak. The shadow battery ---
persistence, smart-persistence, linear lags, size-matched ESN, NVAR, Haar
control, classical-only ablation --- runs on every production window
permanently: champion--challenger is the steady state, not a launch event.

\emph{L4, decision layer.} Calibration wrapper, stacker over the hybrid
feature block, regime router, and cost-sensitive thresholds per persona
cost matrix, emitting alert objects that carry their own provenance and
their delta over persistence:

\begin{strip}
\begin{lstlisting}[language={},caption={Alert object: honesty as payload.},label={lst:prod-alert}]
{ "site": "...", "event": "ramp_down", "horizon_min": 60,
  "p_event": 0.83, "magnitude_q10_q50_q90": [-0.31, -0.22, -0.12],
  "lead_time_min": 47, "regime": "convective_afternoon",
  "conformal_coverage_target": 0.9, "run_id": "2026-07-11T04:20Z#a1b2",
  "model_version": "hybrid-1.3.0",
  "baseline_delta": {"vs_persistence_F1": +0.11} }
\end{lstlisting}
\end{strip}

\emph{L5, surfaces.} REST and streaming API; a dashboard whose benchmark
table is non-removable by design; alert channels; the agent assistant of
Section~\ref{sec:prod-survey-agent}.

\emph{L6, governance.} Input-drift and rolling-coverage monitors; feature
effective rank as a QRC-specific health metric (rank collapse is the
leading indicator that quantum features have degenerated); alert-fatigue
counters; retraining with walk-forward re-validation; promotion between
model versions only through the pre-registered gates.

\section{Implementation Phases 8--12}
\label{sec:prod-phases}

\paragraph{Phase 8 --- data plane.} Registry connectors and one
private-feed adapter; QC service implementing the mask policy with gap
reports; transform service with a frozen-statistics registry keyed by model
version; a light feature store; a walk-forward replay harness driving
everything through the same code path as live inference. Accept:
bit-identical backtest reproduction from raw data with two commands; the
leakage-audit script passes.

\paragraph{Phase 9 --- serving core.} \code{FeatureProvider}
implementations (Sim, Cached; Hardware stub wired through the Stage-5
compiler as CI); quantile and event heads with per-horizon leak; adaptive
conformal wrapper with a rolling-coverage monitor; shadow-battery service;
latency budget of sub-second simulated features at $n\le11$ on a 5-minute
cadence, alerts within 60 seconds of data arrival. Accept: rolling
empirical coverage within three points of target on held-out streams;
battery columns present in every stored prediction row; cache-hit
accounting in logs.

\paragraph{Phase 10 --- decision layer.} Stacker ablation under the
fairness rule (ridge versus gradient boosting versus shallow TCN/LSTM
versus a small transformer, all offered to the classical-only ablation);
regime router (clear-sky classification plus changepoint detection) with
regime-conditional error reporting; a cost-sensitive threshold optimiser
taking the persona cost matrix (expected-cost minimisation, not F$_1$
maximisation); alert engine with deduplication, escalation, and fatigue
caps. Accept: decision-layer value reported in cost units alongside F$_1$;
the stacker table includes the classical-only column, and if the quantum
delta is zero the table says so.

\paragraph{Phase 11 --- surfaces and agent.} API with documentation and
client snippets; dashboard; the assistant of
Section~\ref{sec:prod-survey-agent} with its tool registry
(\code{get\_forecast}, \code{get\_alert\_rationale},
\code{run\_whatif\_backtest}, \code{get\_benchmark\_table},
\code{get\_model\_card}, \code{configure\_site\_draft}), retrieval index
over model cards, run logs, and event post-mortems, the numeric-provenance
guardrail, and the recommend-never-actuate boundary; report generator;
onboarding copilot producing draft site configurations for human approval.
Accept: a grounded-answer evaluation set passes with every numeric claim in
sampled transcripts tracing to a tool call --- zero fabricated numbers
tolerated; time-to-first-alert for a new site under one day.

\paragraph{Phase 12 --- pilot and promotion.} One or two design-partner
sites in shadow mode for at least six weeks against the operator's
incumbent, producing the honest benchmark table; pilot exit criteria
written before the pilot starts (for example, ramp F$_1$ at least the
incumbent's plus $0.05$ at the operator's false-alarm budget, latency SLA
met, zero provenance violations); activation on pass; quarterly promotion
reviews re-running the pre-registered gates, including the contractual
possibility of demoting the quantum feature block. Accept: the pilot
report itself, whatever it says.

\section{Economics and operational honesty}
\label{sec:prod-econ}

The cost axis is reported next to every accuracy number: simulated feature
cost per prediction, cache-hit rate, and --- for any hardware experiment
--- shots and wallclock per prediction; a model needing $10^9$ shots to
match a laptop has answered the business question. Public sources
(SURFRAD, NSRDB, ERA5, NOAA) carry attribution requirements; site SCADA is
private and stays on the customer boundary, with only derived features
leaving if contracted. The degrade-graceful ladder --- quantum features,
then classical-only stacker, then smart-persistence --- keeps the alert
schema identical at every rung, so downgrades are invisible to
integrations.

\section{Survey: classical ML/DL and agentic AI at the post-processing end}
\label{sec:prod-survey}

\subsection{Statistical post-processing}
Raw regression outputs, reservoir or otherwise, are not calibrated decision
inputs. The forecasting tradition supplies the fixes: distributional
regression in the MOS/EMOS lineage, quantile regression on pinball loss,
isotonic recalibration of event probabilities, CRPS as the proper score for
the quantile fan. All are computationally trivial, all sit strictly
downstream of the feature stream, and none care whether the features were
quantum. Adopted: quantile heads plus isotonic recalibration.

\subsection{Conformal prediction: guarantees instead of accuracy claims}
Split and adaptive conformal methods wrap any point or quantile forecaster
and deliver distribution-free finite-sample coverage; adaptive variants
adjust the miscoverage rate online and so track regime shifts. The
strategic fit is exact: the product can contractually promise
``90-percent intervals that cover 90 percent, monitored weekly'' --- a
guarantee --- while accuracy claims remain comparative and scoped.
Adopted: an adaptive conformal wrapper as the default interval layer, with
the rolling-coverage monitor in L6.

\subsection{Machine-learned stackers over hybrid features}
Gradient-boosted trees remain the strongest tabular learners for stacking
heterogeneous blocks: reservoir features, classical lags, calendar and
geometry, NWP and DL-weather channels. The C-QRC preprint is the direct
precedent --- gate-QRC used as a feature extractor feeding LSTM,
Transformer, and XGBoost heads, with the honest finding that gains
concentrate in high-volatility and ramp regimes rather than average point
accuracy. That is this project's thesis stated by an independent group, and
it fixes the stacker's role: event-regime lift, not accuracy everywhere.
Adopted: a gradient-boosted stacker as the default, under the fairness
rule.

\subsection{Deep heads and foundation models: two distinct roles}
Sequence heads (TCN, LSTM, small transformers) are alternative stackers
where horizon-coupled multivariate structure matters, and are evaluated in
the Phase-10 ablation rather than assumed. Time-series foundation models
join the \emph{baseline battery}: if a frozen pretrained forecaster matches
the pipeline on event metrics, the honest table says so. Deep-learning
weather models and classical NWP enter as \emph{exogenous input providers}
--- their gridded forecasts become L2 channels --- because the product's
edge is the event-focused last mile, not global weather modelling. This
division of labour is the survey's central architectural conclusion.

\subsection{Regime-aware routing}
Changepoint and hidden-Markov detectors with clear-sky classification gate
per-regime read-outs and thresholds; regime-conditional error reporting
doubles as the router's evaluation. The Track~A/B architecture routing of
Section~\ref{sec:prod-def} is the same idea one level up: measured memory
requirements select the reservoir.

\subsection{Agentic AI at the user boundary}
\label{sec:prod-survey-agent}
The 2024--2026 literature has converged on language-model agents as the
interpretation and decision-support tier over energy forecasting systems,
not as forecasters. A 2026 review of LLM-agent renewable forecasting
(arXiv:2605.25141) proposes a six-layer taxonomy whose top layers are
uncertainty estimation and natural-language reporting, and names
hallucination control, drift, and EMS interoperability among twelve open
challenges; prototype systems --- GridMind for multi-agent power-system
analysis (arXiv:2509.02494), Grid-Agent for grid control with stepwise
rationales (arXiv:2508.05702), ``Power Agent'' concept papers, and
forecasting-process automation frameworks --- cluster into six patterns, of
which this design adopts five and rejects one:

\begin{enumerate}
\item a conversational analyst over the forecast and metrics store
  (retrieval plus metric tools): questions about yesterday's ramp are
  answered by tool calls whose numbers carry run identifiers ---
  \emph{adopted};
\item an alert triage and explanation agent attaching rationale to each
  alert: regime label, leading feature attributions, retrieved analogue
  past events with outcomes --- \emph{adopted}, and the single
  highest-leverage adoption feature, because operators trust alerts they
  can interrogate;
\item a what-if and scenario agent orchestrating parameterised backtest
  calls (replay a monsoon week under a different alert policy and report
  the cost delta), the decision-theoretic direction of the review
  literature mapped onto the Phase-10 cost matrices --- \emph{adopted};
\item a report and compliance generator drafting persona digests,
  post-event analyses, and re-validation reports from templates and
  tool-fetched numbers --- \emph{adopted};
\item an onboarding copilot translating natural language into a draft site
  configuration for human approval --- \emph{adopted}, the
  easy-to-adopt lever, with time-to-first-alert measured in hours;
\item an autonomous control agent actuating switches or dispatch ---
  \emph{rejected}: the boundary is recommend-never-actuate with a human in
  the loop, and actuation belongs to the customer's EMS and its safety
  case, not to a forecasting vendor's language model.
\end{enumerate}

Guardrails, non-negotiable and recognisably this project's honesty
discipline transposed: a numeric-provenance rule (the agent may state only
numbers returned by tools, rendered with run identifiers, enforced by a
verifier pass and audited with zero tolerance in Phase~11); uncertainty
phrasing bound to the conformal coverage actually monitored; refusal
templates for horizons and regimes outside validated scope; the benchmark
table retrievable in one tool call, so ``how much better than
persistence?'' is always answerable; and no advantage language --- the
agent describes quantum features as quantum-inspired simulated reservoir
features with the measured ablation delta. Orchestration starts with a
single tool-calling agent; a planner--verifier pair is added only if
transcript audits show the need; tool interfaces follow an open
tool-protocol convention so the assistant can later live inside the
customer's own chat tools.

\subsection{Adoption design}
An integration ladder in which each rung is a complete product: alerts by
mail or chat with zero integration; webhook, CSV, and REST pull; EMS and
SCADA integration through partner historians with adapters on the customer
side; finally the assistant embedded in the operator's chat tools. Trust
artefacts shortcut sales: a six-week free shadow mode against the
operator's incumbent ending in the honest benchmark table --- the
no-overclaiming discipline as the sales instrument; the open benchmark
harness, so claims are independently re-runnable; and model cards carrying
the pre-registered criteria and their outcomes, including nulls. Pricing is
sketched per site, tiered by alert budget and horizon set, with a
researcher tier at cost.

\section{Risks specific to the product layer}
\label{sec:prod-risks}

\begin{table*}[h]
\centering\small
\caption{Product-layer risks and responses.}
\label{tab:prod-risks}
\begin{tabular}{@{}p{5.4cm}p{8.6cm}@{}}
\toprule
Risk & Response \\
\midrule
agent states a fabricated number & provenance verifier and zero-tolerance evaluation gate; numbers only via tools \\
alert fatigue kills adoption & cost-matrix thresholds, fatigue caps, weekly false-alarm budgets per persona \\
quantum delta $\approx0$ in production & contractual quantum-optional demotion; value rests on decision layer and governance either way \\
exogenous ramp drivers dominate & NWP and DL-weather channels in L2; sky-imager partnership on the roadmap, not the critical path \\
a foundation-model baseline matches the pipeline & it joins the battery and the table says so; the moat is then event-regime calibration and the decision layer, and the table proves whether that holds \\
data-licence creep on private feeds & derived-features-only export; raw stays on the customer boundary \\
\bottomrule
\end{tabular}
\end{table*}

\section{Deliverables}
\label{sec:prod-deliverables}

Extending Part~X's D1--D6: D7, the data plane with its replay harness
(Phase~8); D8, the serving core with its coverage monitor (Phase~9); D9,
the decision layer with cost-unit evaluation (Phase~10); D10, the surfaces
with a provenance-audited assistant (Phase~11); D11, the pilot report with
pre-written exit criteria (Phase~12). Checkpoint handoffs follow D8 and
D10.
